\section{Discussion}\label{sec:discussion}

\subsection{What the theorem says about critical collapse}

In the standard picture of critical collapse~\cite{GundlachMartinGarcia2007}, initial data near the threshold of
black hole formation evolve for a long time close to the critical solution, and the deviation is dominated by its
single growing mode, $e^{\lambda T}$. If the initial data depend on a parameter $p$, the black hole mass then
scales like $(p-p_*)^\gamma$ with $\gamma=1/\lambda$, modulated by a periodic wiggle in the discretely self-similar
case. The Main Theorem supplies the spectral input of this picture for spherically symmetric scalar field
collapse: there is exactly one unstable mode, it is simple, and $\lambda\in[2.674040,\,2.674115]$. Turning the
picture into a theorem needs further ingredients, among them the following; the regularity question of
\cref{sec:discussion-regularity} and the neutral modes of \cref{rem:neutral} would also have to be addressed.
\begin{enumerate}[label=(\alph*),leftmargin=*]
\item \emph{Linear stability in the evolution sense.} Absence of other roots of $L$ in $\Re s>0$ does not by
  itself give decay of the linearized flow on a complement of the unstable and gauge modes: the operator is not
  self-adjoint, and one needs resolvent bounds along vertical lines, or a completeness statement for the Floquet
  modes, rather than the location of eigenvalues alone. Reiterer has given an abstract framework for finite
  codimension stability of time-periodic symmetric hyperbolic equations with a compact resolvent sufficiently far to the right in the complex
  plane, motivated by discretely self-similar spacetimes~\cite{Reiterer2015}; whether its hypotheses hold for the
  present formulation is not addressed here.
\item \emph{Nonlinear stability:} a codimension-one stable manifold of the critical solution.
\item \emph{Global aspects.} The solution of~\cite{RT2019} is constructed on a neighbourhood of the past light cone
  of the singularity; the threshold picture concerns asymptotically flat data.
\end{enumerate}
The numerical agreement of $1/\lambda$ with the measured exponent $\gamma\approx0.374$ is therefore a consistency
check, not a consequence.

\subsection{Nonspherical perturbations}

Reiterer and Trubowitz describe the numerical evidence on general perturbations as inconclusive~\cite{RT2019},
citing the linear analysis of Mart\'in-Garc\'ia and Gundlach~\cite{MartinGarciaGundlach1999}, who found all
nonspherical perturbations to decay, and the axisymmetric nonlinear evolutions
of~\cite{ChoptuikHirschmannLieblingPretorius2003}. The delicate sector is the polar (even) sector with $\ell=2$,
whose least damped mode was found at $\kappa\Delta\approx-0.07$, close to the imaginary axis.

As a first step toward extending the Main Theorem, we have written the Gerlach--Sengupta
equations~\cite{GerlachSengupta1979} on the background of~\cite{RT2019}, in the variables $\omega$ and $\mu$, and
solved them numerically by spectral collocation, keeping the light cone inside the domain and filtering spurious
solutions by regularity and by the constraints. This computation is \emph{not} rigorous and is not part of the
theorem. In both sectors, for $\ell=2$ and $\ell=3$, it is consistent with the values
of~\cite{MartinGarciaGundlach1999} within the precision of their finite-difference computation; for instance
$\kappa\Delta\approx-2.320$ against their $-2.30$ for the axial mode with $\ell=2$. For the polar mode with $\ell=2$
it gives
\[
  s\approx-0.004764\pm0.147345\,i,\qquad \kappa\Delta=4\pi\Re s\approx-0.0599,
\]
against their $-0.07$, and stable under refinement of our discretization to about $2\cdot10^{-5}$, which is
numerical evidence and not a certified bound: the mode is damped, but at a distance of only $0.005$ from the
axis. A rigorous count for $\ell\ge1$ appears feasible with the methods of this paper, using a deflation and a
Newton--Kantorovich enclosure for this mode, the constraint structure in the certificate, and a bound uniform in
large $\ell$. Together with the Main Theorem it would give exactly one unstable mode for general linear perturbations,
in the spectral sense.

\subsection{Regularity}\label{sec:discussion-regularity}

The class of physical modes in \cref{def:physical-mode} requires analyticity in $\xi$ across the light cone and
at the axis, with a complex neighbourhood that may be arbitrarily thin. Whether every Floquet mode with
$\Re s>0$ that is merely smooth in $\xi$ is analytic is an open regularity question for a hyperbolic system that is
periodic in $\tau$. Our radius recovery (\cref{lem:recovery}) works for every weight $\kappa_2'>1$, but its bound on
the free part degenerates like $1/(\kappa_2'-1)$, so it does not reach $\kappa_2'=1$.

\subsection{On the method}

What made the count feasible on ordinary computers was a combination of choices of formulation.
\begin{itemize}[leftmargin=*]
\item \emph{A finite reference instead of a majorant of the exterior.} An earlier approach bounded the resolvent of
  the finite section and the operator outside it separately; it required truncations and constants out of reach
  without high-performance computing. Comparing $L$ with a block-diagonal reference whose zeros are those of a
  finite matrix needs only bounds along the contour.
\item \emph{A space in which the background is pointwise.} On the torus space the numerical range of the background
  reduces to the symbol, which gave the exclusion bound $2.93$, against $414$ from block-by-block majorants.
\item \emph{The Perron criterion} for block matrices of norms, instead of the spectral norm of such a matrix.
\item \emph{Deflation of the gauge roots,} which lie on or near the contour.
\item \emph{Constraints through $K$.} In a floating-point diagnostic on the truncation $12\times36$, the maximum
  over the boundary of $[1/8,1]\times[-1/4,1/4]$ of the weighted norm of $J(\,\cdot\,+s)^{-1}$ was about $1800$ for
  $L$ and $14$ for $K$, each in the space where its certificate is formulated. The physical classification is therefore obtained by
  proving $K$ injective, rather than by bounding the constraints on root spaces of $L$.
\end{itemize}
Several other routes were tried and closed, with measured lower bounds on their cost. They are documented in the
repository, for readers who might attempt similar computations.
